\documentclass[11pt]{article}
\usepackage[margin=1in]{geometry}
\usepackage{amsmath,amssymb,mathtools,mathrsfs,url}
\setlength{\parindent}{0pt}
\setlength{\parskip}{5pt}
\setlength{\emergencystretch}{2em}
\DeclareMathOperator{\Tr}{Tr}
\DeclareMathOperator{\tr}{tr}
\DeclareMathOperator{\diag}{diag}
\begin{document}
\paragraph{Nonlinear causal state, clock elimination and a finite-parent test
[X/S; V:X,W,D].}
This result specifies an auxiliary mean-state time prescription for the
unchanged sine-polar $H_*$.  It proves a local nonlinear construction using
its exact clock allocations, and gives one bounded implementation
specification for a declared $g_M=g_B=0$ coupling control.  The theorem
uses the complete interacting light law at fixed finite couplings; the
Gaussian implementation does not.  Neither construction minimizes the
sea at every time, truncates a moving state to its pair sector, fixes
$A=B$, or prescribes an inward characteristic.  The single independent
review \path{solve/review_X_nonlinear_response.tex} passes this bounded
result and its Gaussian implementation specification after the
incorporated repairs.  Exactly two claim-specific [X] groups apply:
the reduced state generator and conserved collar Hamiltonian, and
the strict clamped flat clock coercivity with its kernel.  The
remaining checked identities and bounds have supporting [S] scope,
with the named methods in SOURCES.  No numerical result is asserted.

\paragraph{1. Exact finite-parent functional and domain.}
Let $\Lambda=I\sqcup B$ be a finite native parent, with all coefficients
crossing the interior/collar partition retained.  Set $\psi_B=N_B=1$;
$T$ is this collar clock.  The state is expressed in one fixed native
CAR Hilbert space.  Use its full finite fermionic Fock space and the
exactly invariant empty sector of the untruncated fixed-CCR bosons.
The latter retain their original vacuum scalar.  This is not a
boson-number modification of the law.  Other boson states require
additional invariant-domain and moment estimates, and are outside the
local finite-dimensional theorem below.

Write the complete lapse-affine expectation as
\begin{equation}
 {\cal H}(N,\psi,\rho)=\sum_{x\in\Lambda}N_x C_x(\psi,\rho),
 \qquad
 C_x=w_x(\psi)+\Tr\{\rho\,e_x(\psi)\},\qquad e_x\geq0 ,
 \tag{XN1}
\end{equation}
where $\rho\geq0$, $\Tr\rho=1$, and
\[
 w_x=\psi_x\tr_{\rm spin}
       \sum_{s=1}^4\eta_s(E_{s1}-E_{s0})_{xx},\quad
 A_s=m_s\Psi^{-1}+i\Psi^{-3}D,\quad
 E_{s0}=|A_s|,\quad E_{s1}=|A_s^\dagger|.
\]
Here $\eta=(+,-,-,+)$ and the squared masses are
$(x_0,x_0+h_1,x_0+h_2,x_0+h_1+h_2)$ with $x_0,h_1,h_2>0$.
The $e_x$ are the original complete positive squares in native
coordinates.  Their light component contains the actual pp/ph/hh
seed, moving excitation filter, four-$B_1$ denominator and the full
$n-1$ middle factor, as in W5--W15.  The splitting, seed and completion
are independent of $N$.  All field derivatives hold the native state
fixed.  On a compact positive ruler chart the polar roots have a gap.
Require also the actual interpolation and ruler-distance prescription
in the seed to be smooth on that chart; a shortest-path minimum with
unresolved ties is not made smooth merely by a polar gap.

The exact vacuum functional is used in (XN1).  Its decomposition into
$W_2=-k_*\int\nabla\psi\cdot\nabla(N\psi)$ and the complete remainder
does not remove any term from the construction.  The only match is
$k_*=k_a(1)=(2\pi G)^{-1}$.  W's continuum coefficient, moving-filter
identity and boundary work retain their separate assumptions.

\paragraph{2. Eliminate the clock, not the causal state.}
Define the rectangular fixed-state derivative
$J_{xy}=\partial_{\psi_y}C_x$, $x\in\Lambda$, $y\in I$.
At a solution $C_I=0$ with invertible $J_{II}$, the implicit function
theorem defines $\psi_I(\rho)$.  The complete ruler equation then gives
\begin{equation}
 J_{II}^{T}N_I+J_{BI}^{T}{\bf1}_B=0,\qquad
 N_I=-J_{II}^{-T}J_{BI}^{T}{\bf1}_B .
 \tag{XN2}
\end{equation}
This is a full transpose solve, including collar source derivatives.
It does not replace those derivatives by a prescribed boundary
normalization.  Positivity of its solution is an additional branch
property.

The reduced autonomous generator is
\begin{equation}
 {\cal E}(\rho)=\sum_{b\in B}C_b(\psi(\rho),\rho),\qquad
 \delta{\cal E}
   =\Tr\{H_{\rm exc}^{\rm nat}(N(\rho),\psi(\rho))\,\delta\rho\},
 \quad H_{\rm exc}^{\rm nat}=\sum_xN_xe_x .
 \tag{XN3}
\end{equation}
Indeed $J_{II}\delta\psi_I=-(\Tr e_x\delta\rho)_{x\in I}$.
Substitute this into $\delta\sum_b C_b$ and use (XN2); the coefficients
of the interior state rows become $N_I$.  Equivalently, restrict the
native state action to $C_I=0$.  Its symplectic term has no ruler
dependence, and the ruler equation cancels the field chain rule.
Both routes give
\begin{equation}
 i\dot\rho=[H_{\rm exc}^{\rm nat}(N(\rho),\psi(\rho)),\rho],
 \qquad \dot{\cal E}=0 .
 \tag{XN4}
\end{equation}
The scalar vacuum is still present in the fields and in ${\cal E}$;
adding it to the state generator changes only the representative phase.
Finite-dimensional local existence and uniqueness follow from the
smooth implicit branch and the ordinary ODE theorem.  The commutator
preserves the density-matrix spectrum, positivity, trace, and every
fixed native charge commuting with the full Hamiltonian.

There is an exact positive charge identity on this finite constraint
surface:
\begin{equation}
 {\cal E}=\sum_{x\in\Lambda}
       \frac{\Tr(\rho e_x)}{\psi_x}\geq0.
 \tag{XN5}
\end{equation}
To see it, divide $C_x$ by $\psi_x$, sum, and use the finite polar
isospectral trace $\sum_xw_x/\psi_x=0$ and $\psi_B=1$.
This identity does not identify a finite box with an isolated
$1/r$ exterior.  At an approximate root the discrepancy is exactly
$-\sum_{i\in I}C_i/\psi_i$ when ${\cal E}$ denotes the collar sum.

\paragraph{3. The actual flat clock branch exists.}
Let $\rho_0=|\Omega_1\rangle\langle\Omega_1|$ be the complete protected
flat reference.  At $\psi=N=1$, $C_x=0$.  Positivity and the common
kernel imply
$\langle\Omega_1,\partial_{\psi_y}e_x\,\Omega_1\rangle=0$.
Thus $J$ at this point is precisely the original vacuum Jacobian,
at unchanged finite couplings.

Here is a finite-polar coercivity proof which needs no gradient
approximation.  For a real site-scalar $\Phi=\diag(\phi)\otimes I_2$
and $E_s=(m_s^2+D^2)^{1/2}$, let
$\mathscr L_s(Y)=E_sY+YE_s$.  Direct root differentiation gives
\begin{equation}
 \delta(E_{s1}-E_{s0})
  =\mathscr L_s^{-1}
          \{-2im_s[D,\Phi]-3[D,[D,\Phi]]\}.
 \tag{XN6}
\end{equation}
The extra clock factor $\psi$ contributes zero since the flat root
difference is zero.  In a $D$ eigenbasis, exchanging its two indices
cancels the first commutator in the real quadratic form.  Consequently
\begin{equation}
 -\phi^T J_0\phi
   =3\langle[D,\Phi],\mathscr F_D[D,\Phi]\rangle_{\rm HS},
 \qquad \mathscr F_D=\sum_s\eta_s\mathscr L_s^{-1}.
 \tag{XN7}
\end{equation}
For eigenvalues $d_j,d_l$ of $D$, the weight is the mass-rectangle
difference of
\[
 f_{jl}(x)=\frac1{\sqrt{x+d_j^2}+\sqrt{x+d_l^2}}
   =\frac12\int_0^1[x+td_j^2+(1-t)d_l^2]^{-1/2}\,dt .
\]
Its second derivative is a positive integral.  Therefore
\[
 (\mathscr F_D)_{jl}\geq
   w_*:=\frac{3h_1h_2}{8(m_4^2+\|D\|^2)^{5/2}}>0 .
\]
Let
$g_D=\inf_{\phi_B=0,\|\phi_I\|_2=1}\|[D,\Phi]\|_{\rm HS}^2$.
If $B$ meets every connected component of the graph of the nonzero
blocks $D_{xy}$, then $g_D>0$.  Hence
\begin{equation}
 -\phi_I^T J_{0,II}\phi_I\geq
       \gamma_0\|\phi_I\|_2^2,\qquad
 \gamma_0=3w_*g_D,\qquad
 \|J_{0,II}^{-1}\|_2\leq\gamma_0^{-1}.
 \tag{XN8}
\end{equation}
Symmetry of $J_0$ is not needed for this inverse conclusion.
On the Cartesian open sine parent it also follows from D5--D9.
An independent positive heat representation of (XN7) is
\[
 -\phi^TJ_0\phi=\frac3{2\sqrt\pi}\int_0^\infty
 t^{-1/2}e^{-tx_0}(1-e^{-th_1})(1-e^{-th_2})
 \int_0^1
 \left\|e^{-(1-u)tD^2/2}[D,\Phi]e^{-utD^2/2}\right\|_{\rm HS}^2
 \,du\,dt .
\]
It follows by inserting the elementary Laplace integral for the
inverse square root into $f_{jl}$; no continuum spectrum is used.
Every finite heat operator is invertible, so this route gives the
same clamped commutator kernel directly.
The full summed linear vacuum derivative is zero at the flat point:
differentiate $\Tr(E_{s1}-E_{s0})=0$ and use $E_{s1}=E_{s0}$ there.
Thus $J_{0,II}^T{\bf1}_I+J_{0,BI}^T{\bf1}_B=0$ and (XN2) gives $N=1$.

It follows that an open neighborhood of $\rho_0$ has an actual
positive-clock-root and positive-lapse branch.  This is a local
finite-parent nonlinear mean-state theorem at every fixed finite
coupling for which the stated source chart is smooth.  It assumes no
adiabatic sea elimination or full-Fock excitation gap.  Its
neighborhood need not be uniform in volume, a small light mass or
coupling.

For a quantitative sufficient test, put $b=\|J_{0,II}^{-1}\|_2$.
On a radius-$r<1$ ruler ball and a fixed physical state, suppose
$q=b\sup\|J_{II}-J_{0,II}\|_2<1$ and
$b\|C_I(1,\rho)\|_2\leq(1-q)r$.  The map
$\psi\mapsto\psi-J_{0,II}^{-1}C_I(\psi,\rho)$ is a contraction.
The further sufficient estimate
\begin{equation}
 \|N_I-{\bf1}_I\|_2
 \leq\frac b{1-q}
 \bigl\|(J_{\Lambda I}-J_{0,\Lambda I})^T{\bf1}_\Lambda\bigr\|_2<1
 \tag{XN9}
\end{equation}
keeps every lapse positive.  The supremum conditions must be checked
or bounded; a small residual alone does not prove a root.

For an implementation diagnostic at an already evaluated trial field,
put $r^{(1)}=J_{\Lambda I}^{T}{\bf1}_\Lambda$.  The transpose solve
in (XN2) gives, even before the clock residual vanishes,
\[
 J_{II}^{T}(N_I-{f1}_I)=-r^{(1)},\qquad
 \|N_I-{f1}_I\|_\infty
 \leq\frac{\|r^{(1)}\|_2}{\sigma_{\min}(J_{II})}.
\]
In particular, an upper bound strictly below $7/8$ guarantees that
both declared lapse guards pass.  Without an arithmetic enclosure
this is a consistency diagnostic, not a certified neighborhood.
Its converse is false.  Record the interior minimum and maximum
and which guard fires: rejection at $\max N_I\geq8$ is not loss of
lapse positivity, and the collar minimum $N_B=1$ can conceal that
upper excursion.  Failure of a corrector within the declared trial
budget neither disproves the local branch nor places its radius
below the attempted load.  These readbacks use the existing trial
Jacobian and lapse; they introduce no new load, field evaluation,
acceptance rule, or allowance beyond the declared work caps.

\paragraph{4. Tangency and finite-frequency meaning.}
Along (XN4),
\begin{equation}
 J_{II}\dot\psi_I
  =-\bigl(i\Tr\rho[H_{\rm exc}^{\rm nat},e_x]\bigr)_{x\in I}.
 \tag{XN10}
\end{equation}
Differentiating (XN2) determines $\dot N$ with the full derivative of
$J$, not with a frozen vacuum kernel.  Equations (XN2)--(XN4) therefore
provide preservation while the regular branch exists.  Continuing
them requires a compact positive-field region, bounded inverse
Jacobian and the stated source/domain regularity.  Loss of one of
these conditions ends this theorem; it is not a proof of a
long-wavelength trapping surface or of its absence.

In a moving polar chart the same evolution is
$i\dot\rho_{\rm pol}=[H_{\rm exc}-iT_\psi^\dagger\dot T_\psi,
\rho_{\rm pol}]$.  Its complete connection is retained once,
including pair and higher-sector transitions.  For one oriented
packet $({\cal A},E,V)$,
$E(V^\dagger\dot V)+(V^\dagger\dot V)E
=V^\dagger\dot{\cal A}-\dot{\cal A}^\dagger V$.
The native formulation (XN4) needs no added connection term.

The causal sea is already in $\rho(T)$.  There is no new vacuum
kinetic term in (XN1).  The accepted XS frozen-field finite-frequency
kernel is a component of the linearized state/source system, not
by itself the reduced response after the clock and lapse have
been eliminated self-consistently.  Its local adiabatic form
is only an approximation to that component.  In the interacting theorem its
light response uses
$H_{2,\ell}=e^{S_2}{\cal B}_{N,2}e^{S_2}$ and its actual inverse.
An independent scale check must compare the causal response to the
spatial ruler stiffness, not just to a one-excitation mass.  The
finite-parent construction remains meaningful when that derivative
approximation fails, provided its exact branch remains regular.

\paragraph{5. Exact Gaussian specialization for a bounded test.}
The numerical specification selects $g_M=g_B=0$ before outcomes.
This removes the dressing only for this control.  It is not a
finite-coupling Thomas--Fermi calculation or a production parameter
choice.  Keep all four actual fermionic seas and every full native
mode.  For the $+$ orientation take $(E,V)=(E_0,U)$; for the $-$
orientation take $(E,V)=(E_1,U^\dagger)$.  For each packet define
\begin{equation}
 X=E^{1/2},\quad B_N=X(N\Psi)X,\quad
 h=\begin{pmatrix}0&B_NV^\dagger\\VB_N&0\end{pmatrix},\quad
 W_\pm=2^{-1/2}\binom{I}{\pm V},\quad P_-=W_-W_-^\dagger .
 \tag{XN11}
\end{equation}
Let $0\leq C_s\leq I$ be its native one-particle covariance, defined
so that a one-body expectation is $\Tr(h_sC_s)$.
The exact Gaussian functional and its evolution are
\begin{equation}
 \begin{aligned}
 {\cal H}_G&=W_0+\sum_s\{\Tr(h_s C_s)+\Tr B_{N,s}\},&
 i\dot C_s&=[h_s,C_s],\\
 C_x^{G}&=w_x+\sum_s\{\Tr(h_{s,x}C_s)+\Tr B_{s,x}\},&
 R_y^{G}&=\partial_{\psi_y}W_0+
       \sum_s\{\Tr((\partial_{\psi_y}h_s)C_s)
                              +\Tr\partial_{\psi_y}B_{N,s}\}.
 \end{aligned}\tag{XN12}
\end{equation}
Here $h_{s,x}=\partial_{N_x}h_s$ and
$B_{s,x}=X_s(\psi_x\Pi_x)X_s$.
In the instantaneous polar blocks the excitation clock source is
\[
 e_x^G=\sum_s\Tr B_{s,x}
      \{(W_s^\dagger C_s W_s)_{++}
                     +I-(W_s^\dagger C_s W_s)_{--}\}\geq0 .
\]
Cross-band coherence remains in $C_s$.  No sea reset is made.
Gaussianity is preserved exactly under this state-dependent
quadratic Hamiltonian: once its self-consistent field history is
given, each covariance evolves by a one-particle unitary.  Thus
(XN12), its clock root and transpose lapse are a closed exact
nonlinear system at the declared zero couplings, not a nonlinear
pair approximation.

All source derivatives are literal fixed-native derivatives of
(XN12).  Use, for example,
$E\,\delta E+\delta E\,E=\delta({\cal A}^\dagger{\cal A})$,
$X\,\delta X+\delta X\,X=\delta E$,
$\delta V=(\delta{\cal A}-V\delta E)E^{-1}$, with the actual
orientation.  The $W_\pm$ and $\Tr B_N$ derivatives may not be
discarded or counted twice.  Degenerate eigenvalues require
Sylvester divided differences, not derivative choices for
individual eigenvectors.

\paragraph{6. Fixed diagnostic input and preparation.}
Use units $\hbar=c=a=1$ for computation, retaining $a$ in the
definitions.  The selected squared masses are
\[
 a^2(m_1^2,m_2^2,m_3^2,m_4^2)
       =(1/12,1/6,1/6,1/4).
 \tag{XN13}
\]
This is the positive diagnostic rectangle $R_D$ with $aM=1/2$,
not a Tier-1 or physical light hierarchy.  Set the unused seed
parameters $\Lambda=m_1$ and $\ell=a$; at zero couplings they
do not affect any displayed operator.
The Newton readout is only the single finite-spacing match:
use C3's exact full-zone $I(1/2)$ for this rectangle and
$k_*=3I(1/2)/a^2$, $G=(2\pi k_*)^{-1}$.
No value is tuned to a radius, lapse or observed contraction.
The row evolution itself needs no numerical evaluation of $G$.
For the primary/remainder readout, evaluate $I(1/2)$ using C3's
positive grouped heat integral, with relative enclosure width
at most $10^{-8}$ and at most $10^5$ scalar integrand evaluations
inside the same whole-batch time cap.  Retain this scalar
uncertainty in the comparison.  Failure leaves that readout
unresolved and does not replace the exact polar evolution by
a fitted coefficient.

The parent boxes are $L=5$ and $L=7$ sites in each direction,
centered at the origin.  $B$ is the outermost layer and $I$ its
complement.  The original open centered sine derivative has
$(D_i)_{x,x+ae_i}=1/(2ia)$ and the Hermitian reverse entry,
with no periodic wrap and no external sites.  All links between
$I$ and $B$ are retained.  $D=\sum_i\sigma_iD_i$.  There is no
spherical averaging in the equations or exclusion of high native
momenta.  The two boxes test a finite-parent collar dependence;
they do not establish an isolated or continuum limit.

At the flat ruler use each packet's actual $P_{s,-}^{\,0}$.
Let $r=|x|$, $r_0=a$, $\sigma=a/2$, $r_c=7a/4$, and $p_0=1/(4a)$.
Define a single scalar lattice envelope by
\[
 g_x={\cal N}
 \begin{cases}
 \exp\!\left[-\dfrac1{1-r^2/r_c^2}
            -\dfrac{(r-r_0)^2}{2\sigma^2}
            -ip_0\sqrt{r^2+\sigma^2}\right],&r<r_c,\\
 0,&r\geq r_c ,
 \end{cases}
 \qquad \sum_x|g_x|^2=1 .
\]
Use both spin vectors $g\otimes e_\uparrow,g\otimes e_\downarrow$
and embed them with $W_{1,+}^{\,0}$.  Their orthogonal rank-two
projector is $P_g$.  Prepare
\begin{equation}
 C_{1}(0;\lambda)=P_{1,-}^{\,0}+\lambda P_g,\qquad
 C_s(0;\lambda)=P_{s,-}^{\,0}\ (s\ne1),\qquad
 \lambda\in\{1/256,1/16,1\}.
 \tag{XN14}
\end{equation}
This is a valid Gaussian density operator.  Its added mean native
charge is $2\lambda$, with variance $2\lambda(1-\lambda)$.
Only $\lambda=1$ is the pure filled rank-two preparation.
The phase selects an inward seed tendency, not an imposed
trajectory; the actual initial and later radius derivative must
be calculated.  The native state in (XN14) is held fixed while
the initial clock root is continued from $\lambda=0$.
The heavy initial seas consequently need not be the instantaneous
seas at the solved ruler.  This is part of the declared preparation,
and its excitation energy is retained rather than subtracted.

Continue the branch connected to $(\lambda,\psi,N)=(0,1,1)$,
with $\lambda$ steps at most $1/16$, halved after a failed corrector.
Stop this preparation after 96 attempted load-continuation
correctors per box or below a $\lambda$ step $2^{-16}$.
Each such attempt has one field solve, limited to 12 field
evaluation trials, including its initial evaluation, rejected
Newton steps and backtracking trials.  A restart at the same
load counts as another continuation attempt; it does not reset
the counter without cost.  Report every prescribed load
reached, and no substitute load selected for a favorable outcome.
Only reached loads with the regularity conditions below are evolved.
The continuation parameter is a preparation label, not physical time
and not an adjustable coupling.

\paragraph{7. Evolution, budgets and dispositions.}
For each reached load use (XN12) and the exact clock/lapse solves,
up to $T_{\max}=2a$.  Use a symmetric implicit unitary midpoint
iteration: at a step solve the fields at
$\overline C_s=(C_s^n+C_s^{n+1})/2$, and impose
\[
 C_s^{n+1}
  =e^{-i\Delta T\,h_s(\overline C)}C_s^n
                        e^{i\Delta T\,h_s(\overline C)}.
 \tag{XN15}
\]
This specified second-order method preserves covariance eigenvalues
at its converged discrete nodes.  It does not conserve ${\cal E}$
exactly at finite step.  Its charge drift and time convergence are
measured rather than corrected by rescaling a field or energy.
Use the predeclared step ladder
$\Delta T/a=1/8,1/16,1/32$, a maximum of 24 outer covariance
midpoint attempts per step, and no more than 768 accepted time steps in the whole
batch.  Each ladder trajectory starts from the same continued
initial covariances and field root for its box and load, and
uses one fixed declared step.  A failed run may start the next
finer declared run from that initial state; a mixed-step rescue
is not a time-convergence comparison.  No unlisted step or
longer trajectory is automatically authorized.
At each corrector keep all source labels and recompute both the
clock Jacobian and its transpose ruler solve.  Reconstruct these
fields again at every candidate endpoint covariance before
accepting the step.  Endpoint charges, sources and all readouts
use those endpoint solutions, never the midpoint fields.
Require the field residuals and branch guards at both endpoints
and midpoint iterates.  An out-of-domain iterate is not admitted
as a valid midpoint; any backtracking or rejected corrector
still counts toward the declared iteration and whole-batch caps.
Each outer midpoint attempt has one inner clock/lapse field
solve with at most 12 field evaluation trials, including its
initial evaluation and every Newton/backtracking trial.
Restarting that inner solve counts as a new outer attempt.
Endpoint reconstruction has a separate limit of 12 such trials;
exhausting it rejects the step rather than resetting its counter.
The 24 limit is thus an outer-attempt count, not a shared count
of all inner field iterations.  Keep the outer, inner, endpoint
and preparation counters separately in the output.

The initial implementation uses complex binary64 or higher
precision.  Required normalized residuals are $a\|C_I\|_\infty,
a\|R_I\|_\infty\leq10^{-10}$, Hermiticity and covariance spectral
defects at most $10^{-10}$, and nonlinear update residual at most
$10^{-11}$ in Frobenius norm divided by
$\max(1,\|C\|_{\rm F})$.  Time-ladder agreement for each reported
normalized observable must reach $10^{-4}$ between the two finest
completed fixed-step trajectories on the same common reporting
interval.  Both evolve from the shared initial state; partial
histories cannot be joined to create such a comparison.
Require energy drift at most
$10^{-7}\max({\cal E}(0),10^{-6}/a)$.
For the time comparison use
$\|O_1-O_2\|/\max(s_O,\|O_1\|,\|O_2\|)$.
The fixed floors are 1 for fields, speeds and dimensionless
ratios, $a$ for radii and $Z_f$, $a^{-1}$ for individual source
rows and Jacobians, and $\max({\cal E}(0),10^{-6}/a)$ for total
energies.  Matrix comparisons use Frobenius norm; covariances
use floor 1 and the marked material covariance is additionally
divided by its fixed trace $2\lambda$.  A stated operator-norm
or spectral guard is still evaluated in that norm, not replaced
by this reporting convention.
These are declared engineering acceptance tolerances, not
mathematical enclosures or a physical approximation budget.
If directed bounds on accumulated residuals are not supplied,
label the result a converged numerical diagnostic, not a
certified exact trajectory.

For a claimed root enclosure, a Newton residual must be supplemented
by an inverse bound and a local Jacobian-variation bound, as in
(XN9) or a Newton--Kantorovich bound around the actual root.
For trajectory error, a continuous interpolant and an independently
bounded differential defect and local Lipschitz bound permit the
usual Duhamel/Gronwall estimate.  Step comparison alone is not
that estimate.  Source errors must include field errors through
the full fixed-native derivatives.

Predeclare regularity stops: $\min N_I\leq1/8$, $\max N_I\geq8$,
$\min\psi_I\leq1/2$, $\max\psi_I\geq4$, or
$\sigma_{\min}(J_{II})/\sigma_{\min}(J_{0,II})\leq10^{-4}$.
Stop also at a failed required residual, exhausted iteration/step
ladder, nonfinite operator, or resource cap.  These are limits of
this bounded diagnostic, not physical exclusions of a smaller
lapse or another regular branch.  All last valid states and the
actual cause are retained.  Every partial accepted step and
failed corrector evaluation remains charged to the whole-batch
step, evaluation and time caps.  An apparent charge drift above
tolerance first fails the numerical realization of (XN3)--(XN4);
it is not evidence that their exact autonomous charge is drifting.

The whole batch is capped at 6 hours, 24 GiB device memory,
32 GiB host memory, and 16384 complete field/source evaluations.
Every evaluation at a new trial field or covariance is charged
to that last limit, including rejected, backtracking, midpoint,
endpoint and preparation evaluations.  A residual-only trial
still counts even if it reuses rather than rebuilds a Jacobian.
These inner limits add no allowance to the whole-batch caps.
These are upper work limits, not a feasibility promise or
reservation.  Root controls launch admission on the mandated
remote host/environment: free GPU at launch, one own process per
device, nice 10 and no more than 16 own threads.  No execution is
authorized by an analytical formula alone; implementation remains
Sol's task after this single specification review.  Cap exhaustion
is unresolved.  There is no parameter fit or rerun expansion.

\paragraph{8. Required physical and numerical readouts.}
Store the actual native covariances, fields, all source rows and
their derivatives at initial/final states and every $a/8$ of
accepted physical time.  The following quantities belong to the
same state and field solution:
\begin{itemize}
\item The clock and ruler residuals, field positivity, polar gaps,
$\sigma_{\min}(J_{II})$, inverse norm, and the differentiated
clock residual from (XN10).  Verify the finite flat Jacobian
quadratic form (XN7) and its column-sum identity as algebraic
controls before continuation.
\item The collar charge $\sum_b C_b$, full functional
$\sum_xN_xC_x$, positive readout $\sum_x e_x/\psi_x$, and their
residual identities and drifts.  The collar data are fixed, so
there is no imposed boundary work.  Compare the two boxes on
their common interior, including initial state/completion
differences; do not identify box agreement with an isolated limit.
\item For each sea, the actual instantaneous excitation number,
positive clock profile and energy; retain particle--hole
coherence and the full moving polar connection in an independent
representation check.  The native equation has no separate
connection term to add.
\item Propagate the marked material covariance
$M(T)=\lambda U_1(T)P_gU_1(T)^\dagger$ with the same $h_1$.
Its positive density $\tr_{\rm carrier}M_{xx}$ has conserved sum
$2\lambda$.  Use it for the coordinate median, 90-percent
radius and root-mean-square radius, and compute
\[
 \dot r_{\rm rms}
   =\frac{i\Tr\{M[h_1,r^2]\}}
                {4\lambda r_{\rm rms}} .
\]
This marking separates added matter from sea response; it is
not a new clock source.  Neither a thin surface nor a limiting
inward speed is assumed.  Report its collar fraction and the
cubic angular variation.
\item Record $\chi=N\psi$ and $c_{\rm principal}=N\psi^{-2}$
throughout the material quantile shells, their minima and angular
spread.  On these small boxes this is the principal long-wave
coefficient, not an established propagation speed for every
represented wave.  No areal surface is inferred from a cubic
quantile.  A spherical transfer needs resolved radial profiles,
small angular errors, a native momentum window and exterior
convergence before $R=r\psi^2$ and a test of $R=2G{\cal E}$ acquire
that interpretation.
\item Compare the actual scalar rows to the chosen diagnostic
nearest-neighbor primary
$W_{2,a}=-k_*a\sum_{\{x,y\}}(\psi_y-\psi_x)
[(N\psi)_y-(N\psi)_x]$ on every native edge, retaining its collar.
The remainder is exactly $W_0-W_{2,a}$ with the same boundary
allocation.  Report its clock/ruler derivatives and their
relative size wherever the primary denominator is nonzero.
This finite-spacing control does not suppress the remainder.
\end{itemize}

\paragraph{9. Causal-response comparison and the next transfer.}
At each stored state form the actual ruler velocity from (XN10).
Where it is nonzero, normalize its real profile $f$ in Euclidean
site norm.  At that frozen field calculate the free pair response
of each reference sea.  Diagonalize its positive $B_N$ and write
the sea tangent in the corresponding particle--hole basis as
$\eta_{ij}=-(V^\dagger\partial_fV)_{ij}/2$, with the basis
transformation applied to both indices.  Then
\begin{equation}
 Z_f(\omega)=2\sum_{ij}
   \frac{(b_i+b_j)|\eta_{ij}|^2}
        {(b_i+b_j)^2-\omega^2},\qquad
 L_{\rm ad}=\langle\eta_{\dot\psi},H_2^{-1}
                                      \eta_{\dot\psi}\rangle .
 \tag{XN16}
\end{equation}
Here the pair Hamiltonian is free because the selected couplings
are zero, not because a protected one-particle mass is substituted
for an interacting pair gap.  Compute $Z_f(0)$ and the values at
$\omega/\Delta_2=1/8,1/4,1/2$, where $\Delta_2=2\min_i b_i$ over
the reached packets.  Compare $\omega^2Z_f$ to the actual
fixed-lapse vacuum second variation along $f$, and to the
two-derivative stiffness when that comparison has a resolved
spatial meaning.  A zero static denominator is reported as such.
These are response probes, not imposed collapse rates.

The observed history additionally gives a resolved time spectrum
or finite-difference turnover scale of $\dot\psi$; a turnover
ratio is not itself a monochromatic error theorem.  Compare the
actual heavy-sea clock/ruler sources to the XS adiabatic sources,
retaining the specified initial preparation transient.  The
ruler comparison includes
$2\operatorname{Re}\langle\eta_x,\dot\xi\rangle$,
$\xi=H_2^{-1}\eta_{\dot\psi}$, and the XS state/force residual
budget.  Do not certify that force from a first-order state
residual alone.  The occupied light covariance has its actual
in-medium response; a vacuum susceptibility may be reported as
a benchmark, not added to its matter source a second time.
The exact Gaussian trajectory is used whether or not this
adiabatic comparison passes.
Time derivatives beyond (XN10) may use centered differences on
the actual fixed-step histories, with their own step-ladder
comparison; endpoint one-sided values are labelled separately.
Such differences are not analytic derivative enclosures.

This comparison does not require building a fermionic Fock matrix.
In the common polar column basis solve
$B_N\Xi+\Xi B_N=\eta_{\dot\psi}$ and put
$n_\Xi=\Tr(\Xi^\dagger\Xi)$.  The exactly normalized trial vector
$(\Omega+i\xi)/\sqrt{1+n_\Xi}$ has one-particle covariance
\[
 \widetilde C_{\rm ad}=\frac1{1+n_\Xi}
 \begin{pmatrix}
   \Xi\Xi^\dagger&i\Xi\\
   -i\Xi^\dagger&(1+n_\Xi)I-\Xi^\dagger\Xi
 \end{pmatrix}.
\]
Transform it with the actual $W$ before evaluating (XN12).
For any free one-body row with coefficient $O$, its actual/trial
expectation difference is bounded by
$\|O\|\,\|C-C_{\rm ad}\|_1$; scalar row terms cancel in the
difference.  This is an a posteriori source comparison and
does not bound an uncomputed full-state error.
In a locally parallel sea phase the pair-preserving derivative
needed by the acceleration term is the matrix
\[
 {\cal D}_T\Xi=\dot\Xi+\tfrac12[K,\Xi],
 \qquad K=V^\dagger\dot V,\qquad
 \langle\eta_x,\dot\xi\rangle
       =\Tr\{\eta_x^\dagger{\cal D}_T\Xi\}.
\]
The formula uses the common polar basis; a time-dependent
diagonalizing basis must carry its additional connection.
The first-order adiabatic ruler force also differs from the
normalized trial expectation by its explicit normalization
correction.  For real ruler variations
$\langle\eta_x,\eta_{\dot\psi}\rangle
=\Tr(K_x^\dagger K)/4$ is real, so the trial's linear ruler
cross term vanishes.  In the notation of the retained XS
adiabatic expression this gives the explicit comparison
\[
 R^{\rm exc}_{{\rm ad},x}-R^{\rm exc}_{{\rm trial},x}
   =n_\Xi R^{\rm exc}_{{\rm trial},x}
          +2\operatorname{Re}\Tr(\eta_x^\dagger{\cal D}_T\Xi),
 \qquad
 R^{\rm exc}_{{\rm trial},x}
   =\Tr\{(\partial_{\psi_x}h)(C_{\rm ad}-P_-)\}.
\]
The clock comparison similarly has its factor $1+n_\Xi$.
Record the normalization and acceleration terms separately.
This prevents a small trial covariance error from being mistaken
for a complete second-derivative ruler certificate.

The bounded question is whether this same-law causal state
construction supports regular inward motion with a conserved
collar charge, and how much the generated sea source changes
the scalar solution and its principal speed relative to an
instantaneous-reference approximation.  The approximation can
be evaluated as a residual at the exact state/fields; it is not
silently substituted into the evolution.  A root or lapse failure
is located before any proposed continuation.  A quantitative
spherical/open or dressed extension needs, respectively, enlarged
parent/window control or the full W source and many-body state
solver.  Removing the sea connection, re-minimizing it, or using
only a moving pair sector is a named time/state-law change.
No such change is made here.  The physical late-time rate,
mass lock, trapping condition and outgoing massive/thermal
spectrum remain open tests of the resulting self-consistent
history, not inputs to this diagnostic.

\paragraph{Scope and SOURCES delivery.}
The reviewed bounded result is the exact finite-parent local
mean-state reduction with positive
collar charge and the explicit Gaussian zero-coupling control
above.  It constructs the causal state rather than imposing a
vacuum-response kinetic term.  The microscopic full-coupling
formula is retained in the theorem, but only the declared
zero-coupling covariance system is specified for computation.
No isolated/spherical/continuum, physical H2, production-parameter,
long-wavelength trapping, late-time or thermal conclusion follows
without the stated transfer checks.

The decisive calculations are (XN2)--(XN10) in this source.
For the first [X] group, differentiating the implicit clock root
and the full transpose lapse equation gives the same reduced
state generator and conserved collar Hamiltonian as restricting
the native state action.  For the second [X] group, the
spectral Sylvester mass-rectangle weights and the positive
root-heat quadratic form give the same strict clamped sign
and graph kernel.  These are the two independently agreeing
method pairs recorded in
\path{solve/review_X_nonlinear_response.tex}.
The positive charge trace identity, contraction/lapse estimates,
local ODE and tangency consequences, Gaussian normalization and
marked matter, and response/trial-covariance formulas are checked
supporting [S] deductions.  Auxiliary time prescription, diagnostic
inputs and budgets are choices; physical, isolated, stellar and
infinite-domain applications remain [V:X,W,D].  There is no
new [R] or [N] claim.
The finite polar row and its Cartesian coercivity input are
D1--D9 of \path{solve/strand_D_existence.tex}, with
\path{solve/review_D.tex}; the general quadratic-form derivation
is also written here.  The complete differentiated light source
is W5--W15 of \path{solve/strand_W.tex}, with
\path{solve/review_W.tex}.  Exact causal/adiabatic distinctions
and force budgets are XS1--XS18 of
\path{solve/strand_X_causal_sea.tex}, with
\path{solve/review_X_causal_sea.tex}.
The primary shell and charge/lock diagnostics are
\path{solve/strand_X.tex}, with \path{solve/review_X.tex}.
The exact single-match coefficient and diagnostic mass convention
are C1--C2 of \path{solve/strand_C3.tex}; no secondary-branch
numerical value or production coupling is used.
The original free covariance construction underlying (XN11)--(XN12)
is rederived explicitly here from its carrier and polar blocks.
The native-state time prescription, Gaussian preparation, finite
boxes, step ladder and work limits are declared choices.
The independent result disposition belongs to
\path{solve/review_X_nonlinear_response.tex}.  There is no code,
runtime or executed numerical result in this author source.
\end{document}
